\documentclass[a4paper,11pt]{article}
\usepackage[a4paper, top=2cm, bottom=2.5cm, left=2.5cm, right=2.5cm]{geometry}
\usepackage[utf8]{inputenc}
\usepackage[T1]{fontenc}
\usepackage[ngerman]{babel}
\usepackage{microtype}
\usepackage{amsmath,amssymb}
\usepackage{hyperref}
\usepackage{textcomp}
\usepackage{newtxtext,newtxmath}
\usepackage{tikz}
\usepackage{qrcode}
\usepackage{background}

\setlength{\parindent}{0pt}

% Wasserzeichen Einstellungen
\backgroundsetup{
	scale=1,
	color=black!15,
	opacity=0.15,
	angle=45,
	position=current page.center,
	vshift=0pt,
	hshift=0pt,
	contents={%
		\Huge
		Dominic-Ren\'e Schu --- \today --- \url{https://github.com/DominicReneSchu/public}
	}
}

\title{Zusammenfassung der Resonanzfeldtheorie \\ und empirischer Ergebnisse}
\author{Dominic-Ren\'e Schu}
\date{\today}

\begin{document}
	
	\BgThispage
	
	\maketitle
	
	% QR-Code unten rechts
	\begin{tikzpicture}[remember picture,overlay]
		\node[anchor=south east,xshift=-1cm,yshift=1cm]
		at (current page.south east)
		{\qrcode{https://github.com/DominicReneSchu/public}};
	\end{tikzpicture}
	
	\tableofcontents
	\newpage
	
	%% ============================================================
	\section{Einleitung und Zielsetzung}
	%% ============================================================
	
	Die Resonanzfeldtheorie (RFT) ist ein axiomatisches Rahmenwerk,
	das Schwingungskopplung als universelles Organisationsprinzip
	physikalischer und systemischer Ph\"anomene beschreibt.
	Ziel ist eine parametrische Erweiterung der Planck-Relation,
	die Resonanzph\"anomene \"uber Skalen und Dom\"anen hinweg
	einheitlich erfasst -- von Teilchenphysik \"uber Kosmologie
	bis zu Finanzm\"arkten.
	
	Dieses Dokument fasst die theoretischen Grundlagen, die
	Methodik und die empirischen Ergebnisse zusammen. Das
	vollst\"andige Manuskript im IOP-Format ist unter
	\texttt{rft\_manuskript\_de\_iop.tex} verf\"ugbar.
	Alle Herleitungen, Quellcodes und Daten sind offen
	zug\"anglich unter
	\url{https://github.com/DominicReneSchu/public}.
	
	Philosophische und metaphysische Erweiterungen der RFT
	sind in einem separaten Dokument dokumentiert
	(\texttt{RFT\_Philosophie.tex}) und explizit als
	spekulative Perspektiven gekennzeichnet.
	
	%% ============================================================
	\section{Theoretischer Hintergrund}
	%% ============================================================
	
	Die RFT formuliert sieben hierarchisch aufgebaute Axiome
	(A1--A7), aus denen die zentrale
	\textit{Resonanzfeld-Gleichung} folgt:
	\begin{equation}
		E = \pi \cdot \varepsilon(\Delta\varphi) \cdot \hbar \cdot f
	\end{equation}
	Dabei ist:
	\begin{itemize}
		\item $\pi$ die Kreiszahl als Ma\ss{} zyklischer Struktur,
		\item $\hbar$ das reduzierte plancksche Wirkungsquantum,
		\item $f$ die Resonanzfrequenz des Systems,
		\item $\varepsilon(\Delta\varphi) \in [0,1]$ der
		phasenabh\"angige \textbf{Kopplungsoperator}, der die
		St\"arke der Resonanzkopplung zwischen zwei Systemen
		als Funktion ihrer Phasendifferenz $\Delta\varphi$
		beschreibt.
	\end{itemize}
	
	\subsection{Begriffsunterscheidung}
	
	\begin{itemize}
		\item \textbf{Kopplungsoperator}
		$\varepsilon(\Delta\varphi)$:
		Dimensionsloser Parameter in der Grundformel.
		Einfachstes Modell:
		$\varepsilon = \frac{1}{2}(1 + \cos\Delta\varphi)$.
		Grenzwerte: $\varepsilon = 0$ (Antiresonanz,
		destruktive Interferenz),
		$\varepsilon = 1$ (vollst\"andige Phasenkoh\"arenz).
		\item \textbf{Kopplungseffizienz}
		$\eta(\Delta\varphi)$:
		Messgr\"o\ss{}e in den FLRW-Simulationen. Wird aus dem
		zeitgemittelten Kreuzterm zweier gekoppelter Skalarfelder
		extrahiert. Analytische Erwartung:
		$\eta_{\mathrm{theo}} = \cos^2(\Delta\varphi/2)$.
		$\eta$ emergiert als Observable, $\varepsilon$ ist der
		theoretische Operator.
	\end{itemize}
	
	\subsection{Spezialfall Planck}
	
	F\"ur $\varepsilon = 1/(2\pi)$ folgt
	$E = \frac{1}{2}\hbar f$, konsistent mit der
	Grundzustandsenergie des harmonischen Oszillators.
	Die Planck-Relation $E = \hbar\omega$ ist damit ein
	Spezialfall der RFT-Grundformel.
	
	\subsection{Die sieben Axiome}
	
	\begin{enumerate}
		\item[A1] \textbf{Universelle Schwingung:}
		Jede physikalische Entit\"at ist durch eine
		Schwingungsfunktion beschreibbar.
		\item[A2] \textbf{Superposition und Interferenz:}
		Schwingungen \"uberlagern sich linear.
		\item[A3] \textbf{Resonanzbedingung:}
		Kopplung tritt auf, wenn Frequenzen in rationalem
		Verh\"altnis stehen ($f_1/f_2 = n/m$).
		\item[A4] \textbf{Kopplungsenergie:}
		$E = \pi \cdot \varepsilon(\Delta\varphi)
		\cdot \hbar \cdot f$.
		\item[A5] \textbf{Energierichtung:}
		Energie ist eine vektorielle Gr\"o\ss{}e mit Betrag
		und Richtung im Resonanzfeld.
		\item[A6] \textbf{Informationsfluss:}
		Information propagiert ausschlie\ss{}lich entlang
		koh\"arenter Resonanzpfade.
		\item[A7] \textbf{Invarianz:}
		Resonanzstrukturen bleiben unter Variation der
		Beobachterperspektive und Systemparameter erhalten.
	\end{enumerate}
	
	%% ============================================================
	\section{Methodik und Softwarearchitektur}
	%% ============================================================
	
	Die empirische Validierung erfolgt durch drei unabh\"angige
	Simulationspipelines, alle in Python implementiert und
	vollst\"andig im Repository dokumentiert:
	
	\subsection{Monte-Carlo-Resonanzanalyse}
	
	\begin{itemize}
		\item $50\,000$ Pseudo-Experimente auf CMS-Open-Data
		($99\,915$ Dielectron-Ereignisse)
		\item KDE-Hintergrundmodell mit schmalem Signalausschluss
		\item Binomialtest mit Bonferroni-Korrektur \"uber
		68~Fensterbreiten
		\item Bootstrap-Konfidenzintervalle
		($10\,000$~Wiederholungen)
		\item Systematik-Checks \"uber drei KDE-Bandbreiten
		($0{,}3$, $0{,}5$, $0{,}7$~GeV) und zehn unabh\"angige
		Seeds
	\end{itemize}
	
	\subsection{Gekoppelte FLRW-Simulationen}
	
	\begin{itemize}
		\item $1\,530$~Einzelsimulationen
		(51~$H_0$-Punkte $\times$ 30~Phasenwerte)
		\item Klein-Gordon-Gleichung in expandierender Raumzeit
		\item DOP853-Integrator, $t = 0$--$120$,
		$12\,000$~Auswertungspunkte
		\item Jackknife-Fehler auf $d_\eta$,
		Bootstrap-Fehler auf Steigung
	\end{itemize}
	
	\subsection{ResoTrade}
	
	\begin{itemize}
		\item Resonanzfeldtheoretisches Trading-System
		f\"ur den BTC-Markt
		\item AC/DC-Zerlegung (Axiom~1),
		Energierichtungsvektor (Axiom~5),
		Resonanz-Gate (Axiom~6)
		\item 24~Monate Backtest + Live seit Februar 2026
		\"uber Kraken
	\end{itemize}
	
	%% ============================================================
	\section{Ergebnisse}
	%% ============================================================
	
	\subsection{Monte-Carlo auf CMS-Dielectron-Daten}
	
	F\"unf Resonanzen mit empirischem p-Wert~$= 0$ detektiert
	(\"uber $50\,000$ Simulationen, stabil \"uber drei
	Bandbreiten und zehn Seeds):
	
	\begin{center}
		\begin{tabular}{llrrl}
			\hline
			\textbf{Resonanz} & $M_0$ (GeV)
			& \textbf{Hits} & $p_{\mathrm{corr}}$
			& \textbf{emp.\ $p$} \\
			\hline
			$\phi$(1020) & $1{,}020$ & $19\,869$
			& $1{,}07 \times 10^{-54}$ & $0$ \\
			J/$\psi$ & $3{,}100$ & $3\,256$
			& $1{,}02 \times 10^{-121}$ & $0$ \\
			$\Upsilon$(1S) & $9{,}460$ & $4\,186$
			& $3{,}01 \times 10^{-201}$ & $0$ \\
			$\Upsilon$(2S) & $10{,}020$ & $4\,625$
			& $5{,}16 \times 10^{-188}$ & $0$ \\
			Z-Boson & $91{,}200$ & $52$
			& $< 10^{-300}$ & $0$ \\
			\hline
		\end{tabular}
	\end{center}
	
	Best\"atigt Axiom~3 (Resonanzbedingung) und Axiom~7
	(Invarianz \"uber Simulationsdurchl\"aufe).
	
	\subsection{Gekoppelte FLRW-Simulationen}
	
	Zwei skalare Resonanzfelder in expandierender Raumzeit.
	Die Kopplungseffizienz $\eta(\Delta\varphi)$ emergiert
	als Observable und folgt n\"aherungsweise
	$\cos^2(\Delta\varphi/2)$, mit systematischer Abweichung
	durch Hubble-Reibung:
	
	\begin{center}
		\begin{tabular}{lr}
			\hline
			\textbf{Messgr\"o\ss{}e} & \textbf{Wert} \\
			\hline
			Einzelsimulationen & $1\,530$ \\
			Steigung $\mathrm{d}d_\eta / \mathrm{d}H_0$
			& $(0{,}00113 \pm 0{,}00017)$~(km/s/Mpc)$^{-1}$ \\
			$d_\eta$ (Flach, $H_0 = 0$)
			& $0{,}0430 \pm 0{,}0076$ \\
			$d_\eta$ (Planck, $H_0 = 67{,}4$)
			& $0{,}1399 \pm 0{,}0248$ \\
			$d_\eta$ (SH0ES, $H_0 = 73{,}0$)
			& $0{,}1488 \pm 0{,}0263$ \\
			$\Delta d_\eta$ (SH0ES $-$ Planck)
			& $0{,}0063 \pm 0{,}0010$ ($>6\sigma$) \\
			$\Delta\chi^2$ vs.\ Planck-2018-CMB
			& $+16{,}0$ \\
			\hline
		\end{tabular}
	\end{center}
	
	Kontrolltest \"uber f\"unf Expansionsszenarien zeigt
	Anstieg von $d_\eta$ mit $H_0$ im kosmologisch
	relevanten Bereich, mit S\"attigung bei sehr hohen
	Expansionsraten ($\dot{a}_0 > 0{,}6$).
	
	Best\"atigt Axiom~1, 3, 4, 5, 7.
	
	\subsection{ResoTrade}
	
	\begin{center}
		\begin{tabular}{llrl}
			\hline
			\textbf{Marktphase} & \textbf{Zeitraum}
			& \textbf{Trades} & \textbf{vs.\ HODL} \\
			\hline
			Sideways (60k--70k) & M\"ar--Sep 2024
			& 437 & $+33{,}3\%$ \\
			Bullrun (60k--110k) & Sep 2024--M\"ar 2025
			& 429 & $+19{,}8\%$ \\
			Korrektur & M\"ar--Sep 2025
			& 247 & $+4{,}3\%$ \\
			Crash (110k$\to$63k) & Sep 2025--Feb 2026
			& 279 & $+46{,}8\%$ \\
			\hline
			\textbf{Durchschnitt} & \textbf{24 Monate}
			& \textbf{1392} & $\mathbf{+26{,}1\%}$ \\
			\hline
		\end{tabular}
	\end{center}
	
	Best\"atigt Axiom~1, 4, 5, 6.
	
	\subsection{Dynamische Simulationen}
	
	Doppelpendelmodell zur Visualisierung komplexer
	Kopplungszust\"ande:
	\begin{itemize}
		\item Bekannte Bewegungsmuster der klassischen
		Mechanik best\"atigt
		\item Neuartige Resonanzmodi identifiziert
		\item Resonanzpfade und energetische
		Stabilit\"atsr\"aume visualisiert
	\end{itemize}
	
	%% ============================================================
	\section{Mathematische Beweise}
	%% ============================================================
	
	Durch neue Methoden der Reihenentwicklung konnten:
	\begin{itemize}
		\item bisher ungel\"oste Integrale analytisch
		erschlossen werden,
		\item systemische Zusammenh\"ange zwischen Frequenz,
		Kopplung und Energie hergestellt werden,
		\item Br\"ucken zwischen klassischer und
		quantenmechanischer Mathematik geschlagen werden.
	\end{itemize}
	
	%% ============================================================
	\section{Ausblick und offene Fragen}
	%% ============================================================
	
	\begin{enumerate}
		\item \textbf{Unabh\"angige Bestimmung von
			$\varepsilon$:}
		Ab-initio-Herleitung aus der Systemgeometrie.
		
		\item \textbf{Neue Vorhersage:}
		Vorhersage einer bisher unbekannten Resonanz oder
		eines messbaren Effekts, der nur aus der RFT folgt.
		
		\item \textbf{Resonanzreaktor:}
		Resonant gesteuerte Transmutation von Atomm\"ull
		als technologische Anwendung. Die RFT sagt vorher,
		dass durch resonante Anregung bei der Eigenfrequenz
		eines Isotops die Zerfallsrate modulierbar ist --
		im Gegensatz zum Standardmodell, das Kernzerfall
		als rein stochastisch betrachtet. Erste Simulationen
		unter
		\texttt{de/fakten/konzepte/resonanzreaktor/}.
		
		\item \textbf{CMB mit Boltzmann-Solver:}
		Voller $\ell$-Bereich ($\ell = 2$--$2500$) mit
		\textsc{camb} oder \textsc{class} als
		$\Lambda$CDM-Referenz.
		
		\item \textbf{Nichtlineare S\"attigung:}
		Analytische Behandlung der
		$\lambda\varepsilon^4$- und
		$\alpha R\varepsilon^2$-Terme im stark
		expandierenden Regime.
		
		\item \textbf{Live-Validierung:}
		6--12~Monate verifizierbare Live-Performance
		von ResoTrade.
		
		\item \textbf{Weitere Dom\"anen:}
		\"Ubertragbarkeit auf Akustik, Biologie,
		Netzwerktheorie, gekoppelte Oszillatoren.
	\end{enumerate}
	
	%% ============================================================
	\section*{Dokumenten\"ubersicht}
	%% ============================================================
	
	\begin{center}
		\begin{tabular}{ll}
			\hline
			\textbf{Dokument} & \textbf{Inhalt} \\
			\hline
			\texttt{rft\_manuskript\_de\_iop.tex}
			& Physik, Simulationen, Peer Review \\
			\texttt{RFT\_Zusammenfassung.tex}
			& Technische Gesamt\"ubersicht (dieses Dokument) \\
			\texttt{RFT\_Philosophie.tex}
			& Metaphysische Erweiterungen (spekulativ) \\
			\hline
		\end{tabular}
	\end{center}
	
	\section*{Anhang: Verzeichnisstruktur}
	
	\textbf{Monte-Carlo-Simulation:}
	\texttt{de/fakten/empirisch/monte\_carlo\_test/}
	
	\textbf{FLRW-Simulationen:}
	\texttt{de/fakten/simulationen/relativitaet\_verbindung/}
	
	\textbf{ResoTrade:}
	\texttt{de/fakten/empirisch/resotrade/}
	
	\textbf{Resonanzreaktor:}
	\texttt{de/fakten/konzepte/resonanzreaktor/}
	
	\textbf{Dynamische Simulationen:}
	\texttt{de/fakten/simulationen/simulation\_atommodell/}
	
	\vspace{2em}
	
	\textbf{Repository:}
	\url{https://github.com/DominicReneSchu/public}
	\hfill
	\qrcode{https://github.com/DominicReneSchu/public}
	
\end{document}